% Complete Chapter 4 introduction, PDF 43 / printed 33.
% block: hilbert-complexes-heading page:43 type:chapter
\clearpage\setcounter{section}{3}
\section{\BT{Hilbert complexes}{Hilbert 复形}}\label{chap:hilbert-complexes}

% block: hilbert-complexes-introduction page:43 type:ordinary-paragraph
\BT{We have just seen that the differential operators entering the de Rham complex can be viewed as closed unbounded operators between $L^2$ spaces. In this way we obtain a structure which combines the homological algebraic features of a cochain complex, discussed in Chapter~\ref{chap:homological-algebra}, with the functional analytic features of closed unbounded operators on Hilbert space, discussed in Chapter~\ref{chap:unbounded-operators}. This structure is called a \emph{Hilbert complex} and was first promulgated by Brüning and Lesch \BookCite{24}. It is the structure underlying Hodge theory of Riemannian manifolds,}{刚才我们已经看到, de Rham 复形中的微分算子可以视为 $L^2$ 空间之间的闭无界算子. 这样便得到一种结构, 它将 Chapter~\ref{chap:homological-algebra} 中讨论的上链复形的同调代数特征, 与 Chapter~\ref{chap:unbounded-operators} 中讨论的 Hilbert 空间上闭无界算子的泛函分析特征结合起来. 这一结构称为 \emph{Hilbert 复形}, 最早由 Brüning 和 Lesch \BookCite{24} 提出. 它是 Riemannian 流形上 Hodge 理论的基础结构,}%
% block: hilbert-complex-history-footnote page:43 type:footnote
\footnote{\BT{An indication of how natural is the notion of a Hilbert complex is that in the early drafts of \BookCite{13} my coauthors and I introduced this definition, before we discovered that it already occured in the literature \BookCite{24} with the very same name and definition.}{Hilbert 复形概念的自然性, 可以从下面这件事看出: 在 \BookCite{13} 的早期草稿中, 我和合作者引入了这个定义, 此后才发现文献 \BookCite{24} 已经用完全相同的名称和定义提出了这一概念.}}%
% block: hilbert-complexes-introduction-continuation page:43 type:ordinary-paragraph
\BT{and, despite its simplicity, the main results of Hodge theory, for example, the identification of cohomology with harmonic forms and the Hodge decomposition, follow easily in this framework. Associated to each Hilbert complex is a graded map, or sequence of operators, called the abstract Hodge Laplacian of the complex. We explore several different but equivalent formulations of the Hodge Laplacian and demonstrate their well-posedness up to harmonic forms. The case of the de Rham complex on a three-dimensional domain is treated in detail, and the results of the chapter are summarized in Box~\ref{box:hilbert-complexes-summary} at the end of the chapter.}{而且, 尽管这一结构简单, Hodge 理论的主要结论, 例如将上同调与调和形式等同以及 Hodge 分解, 都能在该框架中容易地得到. 每个 Hilbert 复形都伴随一个分次映射或算子序列, 称为该复形的抽象 Hodge Laplacian. 我们将考察 Hodge Laplacian 的几种不同但等价的表述, 并说明在相差调和形式的意义下它们是适定的. 三维区域上的 de Rham 复形情形将得到详细讨论, 本章结果概括于章末的 Box~\ref{box:hilbert-complexes-summary}.}
